Adjudicating between the supply-shock and the monetary-policy accounts of the 1970s turns on timing, on cross-country comparison, and on the estimated behaviour of the central bank's reaction function. \slotmarker{S1} \slotmarker{S2} That recession regularity indicates that oil supply shocks account for the entire rise in US inflation during the 1970s.

The chronology, however, disciplines any such attribution. Only two descriptions of the pre-embargo record are available, and exactly one of them can be correct: either US consumer price inflation exceeded 5 percent before the October 1973 oil embargo, or it remained below 5 percent until that embargo. The published price indices place the decade in the first case, and inflation already above 5 percent in advance of the embargo cannot be squared with the view that oil supply shocks produced the whole of the decade's inflation. What can be affirmed is that at least one of two impulses was in operation, and perhaps both were: OPEC production cuts sharply raised the real price of crude oil in 1973 and 1974, and US money growth accelerated in the years preceding 1973.

The monetary side of the ledger rests on further evidence. Japan experienced lower consumer price inflation than the United States in 1974. The Federal Reserve's policy interest rate stood below the rate of inflation for much of the 1970s; in other words, real policy interest rates in the United States were negative during much of the decade, and the Japanese comparison for 1974 is precisely the sort of international evidence that shows those American real rates to have been negative. That sustained shortfall of real rates below zero is in turn evidence for the estimate of Clarida, Gali and Gertler that the pre-1979 Federal Reserve policy rule violated the Taylor principle: a central bank that raises its nominal rate by less than one-for-one with inflation will let real rates fall as prices accelerate.

The dispute was settled less by any single regression than by an institutional experiment at the close of the decade. \slotmarker{S3} The lesson drawn from that episode is that a determined central bank can end an inflation whose proximate trigger lay in commodity markets, and the weight of professional judgment shifted accordingly toward the monetary account.
